% GENERATED FILE. Edit canonical lessons, then run npm run generate:books.
% canonical-source-hash: fnv1a64:55ad20514c876f7d
% canonical-lessons: KA-C340-dantakathe, KA-C340-purana, KA-C340-daitya, KA-C340-raksasa, KA-C340-sapa, KA-R340-review-334-340

\chapter{Legends and Demons}
\label{ch:ka-legends-and-demons}
\hlchaptermodality{340}

\begin{chapteropening}
\textbf{By the end of this chapter:} \emph{I can say a legend, a myth, a giant, a demon and a curse.}

Complete the last lesson of chapter 340: I can say a legend, a myth, a giant, a demon and a curse.
\end{chapteropening}

\section[\texorpdfstring{\kn{ದಂತಕಥೆ} (dantakathe)}{ದಂತಕಥೆ (dantakathe)}]{\kn{ದಂತಕಥೆ} (dantakathe) — a legend}

\label{lesson:KA-C340-dantakathe}



\begin{quote}
Before the new one: say the Kannada for to scatter, then the Kannada for a stampede.
\end{quote}

\subsection*{You'll want to know: \kn{ದಂತಕಥೆ}}
\textbf{\kn{ದಂತಕಥೆ}} — \emph{dantakathe} — \textquotedblleft{}a legend\textquotedblright{}.

\begin{quote}
\kn{ಈ} \kn{ಕೋಟೆಯ} \kn{ಬಗ್ಗೆ} \kn{ಒಂದು} \kn{ದಂತಕಥೆ} \kn{ಇದೆ}.
\end{quote}

\begin{quote}
\emph{ī kōṭeya bagge ondu dantakathe ide.}
\end{quote}

\begin{quote}
There is a legend about this fort.
\end{quote}

Plural: \kn{ದಂತಕಥೆಗಳು}, \emph{dantakathegaḷu}.

A story told from one generation to the next, partly history and partly not.

\begin{cousinweb}
\textbf{\kn{ದಂತಕಥೆ}} (\emph{dantakathe}) is a Kannada compound whose second part is \kn{ಕಥೆ}, \textquotedblleft{}a story\textquotedblright{}, from Sanskrit \emph{kathā} (CDIAL 2704); the origin of the first part is not settled.
\end{cousinweb}

\begin{grammarlens}[title={What you've built}]
One more word for telling what happened.
\end{grammarlens}

\subsection*{Your turn}
\begin{itemize}
\raggedright
  \item \emph{Say it:} \emph{dantakathe}
  \item \emph{Say it:} \emph{dantakathe}, once more
  \item \emph{Say it:} \emph{ī kōṭeya bagge ondu dantakathe ide.} — There is a legend about this fort.
  \item \emph{Recall:} say the Kannada for to scatter, then the Kannada for a stampede, then say \emph{dantakathe} again
\end{itemize}

\begin{tcolorbox}[breakable,colback=teal!4,colframe=teal!35!black,title={Before you move on}]
What does \kn{ದಂತಕಥೆ} mean? (\textquotedblleft{}A legend\textquotedblright{}.) Say it once more.
\end{tcolorbox}
\section[\texorpdfstring{\kn{ಪುರಾಣ} (purāṇa)}{ಪುರಾಣ (purāṇa)}]{\kn{ಪುರಾಣ} (purāṇa) — a myth; an old legend}

\label{lesson:KA-C340-purana}



\begin{quote}
Before the new one: say the Kannada for a stampede, then the Kannada for a legend.
\end{quote}

\subsection*{You'll want to know: \kn{ಪುರಾಣ}}
\textbf{\kn{ಪುರಾಣ}} — \emph{purāṇa} — \textquotedblleft{}a myth; an old legend\textquotedblright{}.

\begin{quote}
\kn{ಅಜ್ಜಿ} \kn{ಮಲಗುವ} \kn{ಮುಂಚೆ} \kn{ಒಂದು} \kn{ಪುರಾಣ} \kn{ಕಥೆ} \kn{ಹೇಳಿದರು}.
\end{quote}

\begin{quote}
\emph{ajji malaguva muñce ondu purāṇa kathe hēḷidaru.}
\end{quote}

\begin{quote}
Before bed Grandmother told an old legend.
\end{quote}

Plural: \kn{ಪುರಾಣಗಳು}, \emph{purāṇagaḷu}.

The old Sanskrit story collections are the \kn{ಪುರಾಣಗಳು}. In talk, \kn{ಪುರಾಣ} can also be a long, boring tale.

\begin{cousinweb}
\textbf{\kn{ಪುರಾಣ}} (\emph{purāṇa}) is borrowed from Sanskrit \emph{purāṇa}, \textquotedblleft{}ancient, belonging to olden times\textquotedblright{} (CDIAL 8283), then the name of collections of old tales.
\end{cousinweb}

\begin{grammarlens}[title={What you've built}]
One more word for telling what happened.
\end{grammarlens}

\subsection*{Your turn}
\begin{itemize}
\raggedright
  \item \emph{Say it:} \emph{purāṇa}
  \item \emph{Say it:} \emph{purāṇa}, once more
  \item \emph{Say it:} \emph{ajji malaguva muñce ondu purāṇa kathe hēḷidaru.} — Before bed Grandmother told an old legend.
  \item \emph{Recall:} say the Kannada for a stampede, then the Kannada for a legend, then say \emph{purāṇa} again
\end{itemize}

\begin{tcolorbox}[breakable,colback=teal!4,colframe=teal!35!black,title={Before you move on}]
What does \kn{ಪುರಾಣ} mean? (\textquotedblleft{}A myth; an old legend\textquotedblright{}.) Say it once more.
\end{tcolorbox}
\section[\texorpdfstring{\kn{ದೈತ್ಯ} (daitya)}{ದೈತ್ಯ (daitya)}]{\kn{ದೈತ್ಯ} (daitya) — a giant}

\label{lesson:KA-C340-daitya}



\begin{quote}
Before the new one: say the Kannada for a legend, then the Kannada for a myth.
\end{quote}

\subsection*{You'll want to know: \kn{ದೈತ್ಯ}}
\textbf{\kn{ದೈತ್ಯ}} — \emph{daitya} — \textquotedblleft{}a giant\textquotedblright{}.

\begin{quote}
\kn{ಕಥೆಯಲ್ಲಿ} \kn{ಒಬ್ಬ} \kn{ದೈತ್ಯ} \kn{ಬೆಟ್ಟದ} \kn{ಮೇಲೆ} \kn{ಇದ್ದನು}.
\end{quote}

\begin{quote}
\emph{katheyalli obba daitya beṭṭada mēle iddanu.}
\end{quote}

\begin{quote}
In the story a giant lived on the hill.
\end{quote}

Plural: \kn{ದೈತ್ಯರು}, \emph{daityaru}.

As an adjective it means huge: \kn{ದೈತ್ಯ} \kn{ಮರ}, a giant tree.

\begin{cousinweb}
\textbf{\kn{ದೈತ್ಯ}} (\emph{daitya}) is borrowed from Sanskrit \emph{daitya}, the name of a race of demons in the old stories (Monier-Williams); in Kannada it also means a giant.
\end{cousinweb}

\begin{grammarlens}[title={What you've built}]
One more word for telling what happened.
\end{grammarlens}

\subsection*{Your turn}
\begin{itemize}
\raggedright
  \item \emph{Say it:} \emph{daitya}
  \item \emph{Say it:} \emph{daitya}, once more
  \item \emph{Say it:} \emph{katheyalli obba daitya beṭṭada mēle iddanu.} — In the story a giant lived on the hill.
  \item \emph{Recall:} say the Kannada for a legend, then the Kannada for a myth, then say \emph{daitya} again
\end{itemize}

\begin{tcolorbox}[breakable,colback=teal!4,colframe=teal!35!black,title={Before you move on}]
What does \kn{ದೈತ್ಯ} mean? (\textquotedblleft{}A giant\textquotedblright{}.) Say it once more.
\end{tcolorbox}
\section[\texorpdfstring{\kn{ರಾಕ್ಷಸ} (rākṣasa)}{ರಾಕ್ಷಸ (rākṣasa)}]{\kn{ರಾಕ್ಷಸ} (rākṣasa) — a demon; an ogre}

\label{lesson:KA-C340-raksasa}



\begin{quote}
Before the new one: say the Kannada for a myth, then the Kannada for a giant.
\end{quote}

\subsection*{You'll want to know: \kn{ರಾಕ್ಷಸ}}
\textbf{\kn{ರಾಕ್ಷಸ}} — \emph{rākṣasa} — \textquotedblleft{}a demon; an ogre\textquotedblright{}.

\begin{quote}
\kn{ಧೈರ್ಯವಂತ} \kn{ಹುಡುಗ} \kn{ರಾಕ್ಷಸನನ್ನು} \kn{ಸೋಲಿಸಿದನು}.
\end{quote}

\begin{quote}
\emph{dhairyavanta huḍuga rākṣasanannu sōlisidanu.}
\end{quote}

\begin{quote}
The brave boy defeated the ogre.
\end{quote}

Plural: \kn{ರಾಕ್ಷಸರು}, \emph{rākṣasaru}.

A woman demon is a \kn{ರಾಕ್ಷಸಿ}.

\begin{cousinweb}
\textbf{\kn{ರಾಕ್ಷಸ}} (\emph{rākṣasa}) is borrowed from Sanskrit \emph{rākṣasa}, \textquotedblleft{}demoniacal\textquotedblright{}, and as a noun \textquotedblleft{}a demon\textquotedblright{} (CDIAL 10672).
\end{cousinweb}

\begin{grammarlens}[title={What you've built}]
One more word for telling what happened.
\end{grammarlens}

\subsection*{Your turn}
\begin{itemize}
\raggedright
  \item \emph{Say it:} \emph{rākṣasa}
  \item \emph{Say it:} \emph{rākṣasa}, once more
  \item \emph{Say it:} \emph{dhairyavanta huḍuga rākṣasanannu sōlisidanu.} — The brave boy defeated the ogre.
  \item \emph{Recall:} say the Kannada for a myth, then the Kannada for a giant, then say \emph{rākṣasa} again
\end{itemize}

\begin{tcolorbox}[breakable,colback=teal!4,colframe=teal!35!black,title={Before you move on}]
What does \kn{ರಾಕ್ಷಸ} mean? (\textquotedblleft{}A demon; an ogre\textquotedblright{}.) Say it once more.
\end{tcolorbox}
\section[\texorpdfstring{\kn{ಶಾಪ} (śāpa)}{ಶಾಪ (śāpa)}]{\kn{ಶಾಪ} (śāpa) — a curse}

\label{lesson:KA-C340-sapa}



\begin{quote}
Before the new one: say the Kannada for a giant, then the Kannada for a demon.
\end{quote}

\subsection*{You'll want to know: \kn{ಶಾಪ}}
\textbf{\kn{ಶಾಪ}} — \emph{śāpa} — \textquotedblleft{}a curse\textquotedblright{}.

\begin{quote}
\kn{ರಾಕ್ಷಸನ} \kn{ಶಾಪದಿಂದ} \kn{ರಾಣಿ} \kn{ಕಲ್ಲಾದಳು}.
\end{quote}

\begin{quote}
\emph{rākṣasana śāpadinda rāṇi kallādaḷu.}
\end{quote}

\begin{quote}
Through the demon's curse the queen turned to stone.
\end{quote}

A curse is \textquotedblleft{}given\textquotedblright{}: \kn{ಶಾಪ} \kn{ಕೊಡು}, to curse.

\begin{cousinweb}
\textbf{\kn{ಶಾಪ}} (\emph{śāpa}) is borrowed from Sanskrit \emph{śāpa}, \textquotedblleft{}a curse, abuse\textquotedblright{} (CDIAL 12393), from the root \emph{śap}, \textquotedblleft{}to curse\textquotedblright{}.
\end{cousinweb}

\begin{grammarlens}[title={What you've built}]
One more word for telling what happened.
\end{grammarlens}

\subsection*{Your turn}
\begin{itemize}
\raggedright
  \item \emph{Say it:} \emph{śāpa}
  \item \emph{Say it:} \emph{śāpa}, once more
  \item \emph{Say it:} \emph{rākṣasana śāpadinda rāṇi kallādaḷu.} — Through the demon's curse the queen turned to stone.
  \item \emph{Recall:} say the Kannada for a giant, then the Kannada for a demon, then say \emph{śāpa} again
\end{itemize}

\begin{tcolorbox}[breakable,colback=teal!4,colframe=teal!35!black,title={Before you move on}]
What does \kn{ಶಾಪ} mean? (\textquotedblleft{}A curse\textquotedblright{}.) Say it once more.
\end{tcolorbox}
\section[(dialogue)]{Looking back: from A Baby in the Family to Legends and Demons}

\label{lesson:KA-R340-review-334-340}



\begin{quote}
Say the Kannada for a new mother and a birth.
\end{quote}

\subsection*{Your turn}
Say these aloud:

\begin{itemize}
\raggedright
  \item a new mother, a birth, a naming ceremony, a baby shower, a successor
  \item an editor, an article, the front page, a leaflet, a publisher
  \item an army, a commander, a warrior, to fight, to retreat
  \item a passing away, a funeral, a cremation ground, a tribute, to comfort
  \item a rumour, to spread far, speculation, the truth, to listen closely
  \item a riot, an uproar, a crowd, to scatter, a stampede
  \item a legend, a myth, a giant, a demon, a curse
\end{itemize}

\begin{tcolorbox}[breakable,colback=teal!4,colframe=teal!35!black,title={Before you move on}]
A new mother? (\textbf{\kn{ಬಾಣಂತಿ}}, \emph{bāṇanti}.) A birth? (\textbf{\kn{ಜನನ}}, \emph{janana}.) A naming ceremony? (\textbf{\kn{ನಾಮಕರಣ}}, \emph{nāmakaraṇa}.) A baby shower? (\textbf{\kn{ಸೀಮಂತ}}, \emph{sīmanta}.) A successor? (\textbf{\kn{ಉತ್ತರಾಧಿಕಾರಿ}}, \emph{uttarādhikāri}.) An editor? (\textbf{\kn{ಸಂಪಾದಕ}}, \emph{sampādaka}.) An article? (\textbf{\kn{ಲೇಖನ}}, \emph{lēkhana}.) The front page? (\textbf{\kn{ಮುಖಪುಟ}}, \emph{mukhapuṭa}.) A leaflet? (\textbf{\kn{ಕರಪತ್ರ}}, \emph{karapatra}.) A publisher? (\textbf{\kn{ಪ್ರಕಾಶಕ}}, \emph{prakāśaka}.) An army? (\textbf{\kn{ಸೈನ್ಯ}}, \emph{sainya}.) A commander? (\textbf{\kn{ಸೇನಾಪತಿ}}, \emph{sēnāpati}.) A warrior? (\textbf{\kn{ಯೋಧ}}, \emph{yōdha}.) To fight? (\textbf{\kn{ಹೋರಾಡು}}, \emph{hōrāḍu}.) To retreat? (\textbf{\kn{ಹಿಮ್ಮೆಟ್ಟು}}, \emph{himmeṭṭu}.) A passing away? (\textbf{\kn{ನಿಧನ}}, \emph{nidhana}.) A funeral? (\textbf{\kn{ಅಂತ್ಯಕ್ರಿಯೆ}}, \emph{antyakriye}.) A cremation ground? (\textbf{\kn{ಸ್ಮಶಾನ}}, \emph{smaśāna}.) A tribute? (\textbf{\kn{ಶ್ರದ್ಧಾಂಜಲಿ}}, \emph{śraddhāñjali}.) To comfort? (\textbf{\kn{ಸಂತೈಸು}}, \emph{santaisu}.) A rumour? (\textbf{\kn{ವದಂತಿ}}, \emph{vadanti}.) To spread far? (\textbf{\kn{ಪಸರಿಸು}}, \emph{pasarisu}.) Speculation? (\textbf{\kn{ಊಹಾಪೋಹ}}, \emph{ūhāpōha}.) The truth? (\textbf{\kn{ಸತ್ಯ}}, \emph{satya}.) To listen closely? (\textbf{\kn{ಆಲಿಸು}}, \emph{ālisu}.) A riot? (\textbf{\kn{ಗಲಭೆ}}, \emph{galabhe}.) An uproar? (\textbf{\kn{ಕೋಲಾಹಲ}}, \emph{kōlāhala}.) A crowd? (\textbf{\kn{ಜನಸಂದಣಿ}}, \emph{janasandaṇi}.) To scatter? (\textbf{\kn{ಚದುರು}}, \emph{caduru}.) A stampede? (\textbf{\kn{ಕಾಲ್ತುಳಿತ}}, \emph{kāltuḷita}.) A legend? (\textbf{\kn{ದಂತಕಥೆ}}, \emph{dantakathe}.) A myth? (\textbf{\kn{ಪುರಾಣ}}, \emph{purāṇa}.) A giant? (\textbf{\kn{ದೈತ್ಯ}}, \emph{daitya}.) A demon? (\textbf{\kn{ರಾಕ್ಷಸ}}, \emph{rākṣasa}.) A curse? (\textbf{\kn{ಶಾಪ}}, \emph{śāpa}.)
\end{tcolorbox}
